\subsection{Earlier Method and Filter Tests}
\label{sec:eval:ablation}

Table~\ref{tab:repair_ablation} compares four saved settings.
The no-completion setting returns the cleaned graph.
Direct LLM uses its own prompt and skips cleanup.
The filter test scores the same proposal before and after checks.
Section~\ref{sec:eval:mechanisms} tests prompt text.
Section~\ref{sec:eval:individual_ablation} tests each filter and selector step.

\begin{table}[t]
\centering
\caption{Four settings on added-defect inputs. Direct LLM uses a different prompt.}
\label{tab:repair_ablation}
\small
\begin{tabular}{l c c c c}
\toprule
Variant & Repair & Over-repair & F1 & Exact \\
\midrule
No LLM completion & 0.3200 & 0.0000 & 0.8472 & 0.0622 \\
Direct LLM (separate prompt) & 0.9244 & 0.0758 & 0.9520 & 0.8267 \\
No constraint gate & 0.9800 & 0.0227 & 0.9895 & 0.9333 \\
\textbf{Diagnosis + Gate} & \textbf{0.9800} & \textbf{0.0167} & \textbf{0.9917} & \textbf{0.9333} \\
\bottomrule
\end{tabular}
\end{table}

Model completion recovers most missing fields.
Diagnosis + Gate reaches 99.17\% F1, compared with 95.20\% for Direct LLM.
Filtering keeps the same repair rate and reduces over-repair by 0.60 points.
It removes eight triples that fail source matching.
All eight differ from the reference.
